\section{Introduction}\label{sec:introduction}
% Why this review matters (importance/motivation)
% What has been done — identify gaps in the literature
% What is new — your contributions / research questions

Equatorial Plasma Bubbles (EPBs) represent a critical and recurring space weather hazard. Originating from the Rayleigh-Taylor instability in the post-sunset equatorial ionosphere, these large-scale plasma density depletions cause severe amplitude and phase scintillations in trans-ionospheric radio signals. As modern society grows increasingly reliant on Global Navigation Satellite Systems (GNSS) and satellite-based communication networks, EPB-induced disruptions pose a severe threat to global technological infrastructure, necessitating continuous, real-time operational monitoring to mitigate their impact \cite{zhong2025automatic,okoh2025bootstrapping,githio2024machine}.

Historically, the observation and tracking of EPBs have relied on ground- and space-based instruments, primarily All-Sky Imagers (ASIs), Very High Frequency (VHF) radars, and GNSS receivers \cite{zhong2025automatic,githio2024machine,okoh2025bootstrapping}. While these instruments generate vast amounts of high-resolution data, identifying EPB signatures has predominantly depended on manual visual inspection and feature extraction by human experts \cite{okoh2025bootstrapping,zhong2025automatic}. This human-in-the-loop approach is inherently subjective, labor-intensive, and highly susceptible to environmental noise such as cloud cover or lunar illumination \cite{okoh2025bootstrapping}. Furthermore, with the exponential proliferation of space weather observational data---where a single modern imager station can generate terabytes of data annually \cite{zhong2025automatic}---manual processing has become an unsustainable bottleneck, underscoring an urgent need for robust, automated detection frameworks.

To overcome these observational bottlenecks, Machine Learning (ML) and Deep Learning (DL) have emerged as transformative tools in geophysical forecasting \cite{githio2024machine,okoh2025bootstrapping}. By leveraging architectures ranging from traditional classifiers like Support Vector Machines (SVM) and Random Forests to advanced Deep Convolutional Neural Networks (DCNNs), researchers have successfully automated the identification, classification, and morphological segmentation of ionospheric irregularities \cite{githio2024machine,zhong2025automatic,thanakulketsarat2023classification}. These data-driven models not only drastically reduce computational time but also demonstrate a remarkable ability to extract hierarchical features---detecting subtle spatial gradients and complex, non-linear patterns that often elude human observers \cite{okoh2025bootstrapping}.

Despite the rapid adoption of artificial intelligence in space weather monitoring, the literature surrounding ML-based EPB detection remains highly fragmented. Recent studies employ a disparate array of algorithms, data modalities (e.g., raw versus detrended optical imagery, radar plots, and TEC maps), and evaluation metrics \cite{okoh2025bootstrapping}. This lack of standardization complicates direct performance comparisons and obscures the optimal configurations required for real-time operational deployment. Furthermore, recurring challenges---such as model bias due to solar-cycle data imbalances and performance degradation under adverse weather conditions---are frequently noted in individual studies but lack comprehensive, cross-study synthesis.

To address this gap, this study presents a systematic literature review of machine learning applications for the detection of Equatorial Plasma Bubbles. The review synthesizes existing computational approaches, datasets, and evaluation metrics in order to provide a structured overview of the current state-of-the-art. Ultimately, it aims to identify common methodological patterns, highlight recurring challenges, and uncover key gaps that may guide future research efforts.